\documentclass[sigconf]{acmart}

%% Rights info — fill in after rights form submission
\setcopyright{acmlicensed}
\copyrightyear{2026}
\acmYear{2026}
\acmDOI{XXXXXXX.XXXXXXX}

\acmConference[SA '26]{ACM SIGGRAPH Asia 2026 Posters}{December 2026}{Tokyo, Japan}
\acmISBN{XXX-X-XXXX-XXXX-X/26/12}

\begin{document}

\title{Hammerhead HMD: A Cephalophoil-Inspired Underwater Augmented Reality Headset}

\author{Sho Nakazawa}
\affiliation{%
  \institution{University of Tsukuba}
  \city{Tsukuba}
  \country{Japan}
}
\email{shonkzw0927@gmail.com}

%%
%% TODO: 共著者がいれば追加
%%

\begin{abstract}
We present Hammerhead HMD, an underwater augmented reality head-mounted display
inspired by the cephalophoil morphology of hammerhead sharks.
The lateral extension of the housing improves stereoscopic depth perception
and widens the horizontal field of view in subaquatic environments.
We evaluate the hydrodynamic characteristics of the device using
Particle Image Velocimetry (PIV) in a water tunnel,
providing quantitative flow visualization around the headset.
%% TODO: 定量的な結果（抗力係数など）が出たら1文追加
\end{abstract}

\keywords{underwater augmented reality, head-mounted display, hydrodynamics,
particle image velocimetry, wearable}

\maketitle

%%------------------------------------------------------------
\section{Introduction}
%%------------------------------------------------------------

Augmented reality in underwater environments demands optics that withstand
pressure, a housing that resists hydrodynamic drag, and a form factor
compatible with diving equipment.
Conventional terrestrial HMDs are ill-suited for this context:
their boxy profiles generate turbulence, and their narrow baselines
reduce depth cue reliability in water, where light refracts differently
than in air.

We take design inspiration from the hammerhead shark (\textit{Sphyrna} spp.),
whose cephalophoil—the distinctive lateral extension of the head—spatially
separates the eyes to enhance binocular parallax and electroreceptor coverage.
Adopting this geometry for an HMD places display and camera modules at a wider
baseline, improving stereo depth estimation while reducing frontal cross-section.

%%------------------------------------------------------------
\section{System Design}
%%------------------------------------------------------------

%% TODO: システム写真 or CADレンダリングを figures/system.png に追加

\subsection{Structural Design}
The housing extends laterally \textbf{[TODO: 寸法 mm]} on each side from the
diver's midline, forming the cephalophoil profile.
The body is fabricated from \textbf{[TODO: 素材（例: 3Dプリント PLA / PETG / アルミ）]}
and is sealed to \textbf{[TODO: 耐水圧仕様 (m)]} depth.

\subsection{Optical \& Sensing Configuration}
Each lateral pod houses a \textbf{[TODO: カメラ仕様（例: 広角レンズ, 解像度）]}
and a \textbf{[TODO: ディスプレイ（例: micro-OLED, FoV）]}.
Stereo baseline is \textbf{[TODO: mm]}, compared to approximately
65\,mm in a typical terrestrial HMD.

%%------------------------------------------------------------
\section{Hydrodynamic Evaluation}
%%------------------------------------------------------------

\subsection{Experimental Setup}

We evaluated flow behavior using Particle Image Velocimetry (PIV)
in a recirculating water tunnel.
A high-speed camera captured tracer particle images at \textbf{4{,}000\,fps}.
The headset was oscillated at \textbf{5\,Hz} with a stroke of \textbf{9\,mm}
over \textbf{30 steps} to simulate diver head motion during swimming.
PIV analysis was performed using PIVlab~\cite{thielicke2021pivlab}.

%% TODO: 水槽の流速（自由流速 U_inf）を記載
%% TODO: Re数を計算して記載

\subsection{Results}

Figure~\ref{fig:piv} shows \textbf{[TODO: 代表フレーム or 平均流速場の図]}.
The cephalophoil profile produced \textbf{[TODO: 観察された流れの特徴（例: 低乱流域, 後流の対称性など）]}.
%% TODO: 定量結果（例: 最大流速, 抗力推定, 比較対象との差異）

\begin{figure}[h]
  \centering
  %% TODO: PIVlab_out.mp4 から代表フレームを切り出して figures/piv_result.png に保存
  \includegraphics[width=\linewidth]{figures/piv_result}
  \caption{PIV flow visualization around the Hammerhead HMD.
           Vectors show in-plane velocity; color encodes speed magnitude.}
  \label{fig:piv}
\end{figure}

%%------------------------------------------------------------
\section{Conclusion}
%%------------------------------------------------------------

We presented the Hammerhead HMD, an underwater AR headset whose
cephalophoil geometry simultaneously widens stereo baseline and
reduces hydrodynamic resistance.
PIV evaluation confirms \textbf{[TODO: 主な結論]}.
Future work will address \textbf{[TODO: 今後の課題（例: 実使用時の映像評価, 深度推定精度の検証）]}.

%%------------------------------------------------------------
\bibliographystyle{ACM-Reference-Format}
\bibliography{references}

\end{document}
